\section{Introduction}
\label{sec:introduction}

Battery energy storage systems in distribution-network microgrids support temporal arbitrage among renewable generation, load demand, and electricity prices through sequential dispatch decisions. Each charge or discharge command propagates through state-of-charge dynamics, electrochemical losses, network constraints, and future operational flexibility\cite{nakabi2021segan,zhang2023multiagentECM,monfaredi2023multiagent}. Classical optimization remains effective when models and forecasts are accurate. Stochastic renewables, flexible demand, and distributed storage have nevertheless motivated learning-based controllers that adapt from operational data\cite{wang2023collaborative,lin2024coordinated,li2024deep,raj2024reinforcement,li2025optimizing,omar2025hierarchical,zhang2026safeTCE}. Deep reinforcement learning (DRL) has therefore become a relevant tool for microgrid energy management, with reported applications in multi-agent coordination\cite{XU2023106774,zhang2024energy}, degradation-aware scheduling\cite{li2025battery,park2026worldmodel}, and resilient operation\cite{delavari2024rlRobust,momen2024resilience,amiri2026explainableECM}. Recent EPSR studies also use DRL for microgrid economic dispatch, active distribution-network scheduling, photovoltaic-driven reconfiguration, and explainable voltage regulation under renewable-rich operating conditions\cite{liu2025epsrInfoDRL,xue2024epsrSafeDRL,wu2025epsrHierarchical,zhang2025epsrExplainable}.

Evaluating DRL algorithms as candidate operational controllers exposes a validation gap: \emph{cumulative economic reward does not uniquely characterize operational behavior}. A low-cost trajectory may park the battery near state-of-charge boundaries, suppress useful cycling, or satisfy constraints mainly through external safety-filter corrections rather than through the learned policy itself. In distribution networks, where battery dispatch affects voltage profiles, loading limits, and reliability, controller selection requires trajectory-level operational assessment in addition to scalar performance metrics---an aspect not captured by aggregate-return reporting alone\cite{henderson2018deepRLmatters,agarwal2021statistical}.

This challenge is closely related to safe reinforcement learning. Battery storage operates under hard physical constraints, including state-of-charge bounds, power limits, and network feasibility, as well as soft operational preferences such as mid-band inventory maintenance, price-responsive behavior, and terminal convergence. The safe DRL literature has developed constrained Markov decision process formulations\cite{achiam2017cpo}, control-barrier certificates\cite{brunke2022safeLearning}, model-predictive safety layers\cite{odriozola2023shielded}, and action-shielding mechanisms\cite{zhou2023safeRegion,mazumdar2024stopping,misra2024psafe,wang2025shieldedPlanning}. Power-system studies cover Volt--VAR control\cite{wang2020safeVoltVar}, distribution operation\cite{li2022safeDistribution}, voltage support\cite{bi2024safeVoltage}, and inverter regulation\cite{shuai2024safeInverter,omar2025interiorPoint,liu2026digitalTwin}. These safety mechanisms also create a reporting issue: \emph{when low evaluated cost is obtained mainly through shield corrections, the measured outcome can overstate autonomous policy capability}. This distinction affects engineering trust, failure-mode analysis, and qualification.

For battery storage, the distinction is important because the state-of-charge trajectory is itself an operational observable. A policy may record low cost while dwelling at inventory extremes, missing terminal targets, maintaining insufficient discharge reserve during high-price periods, or producing charge--discharge patterns inconsistent with intertemporal value. Such behaviors are not fully captured by feasibility flags or cumulative metrics. We define \emph{battery morphology} as the set of trajectory-level attributes---mid-band residence, boundary avoidance, terminal consistency, price-responsive directionality, infeasible-action suppression, network constraint satisfaction, and intervention diagnostics---used to assess whether a learned policy produces interpretable energy shifting, rather than only low evaluated cost after protective intervention.

This perspective is consistent with the broader view that algorithmic evaluation extends beyond aggregate performance metrics\cite{henderson2018deepRLmatters,agarwal2021statistical}. Studies of DRL reproducibility, transfer learning\cite{padakandla2021dynamicRL,zhu2023transferRL}, verification\cite{landers2023verification}, and explainability\cite{vouros2022xrl,milani2024xrl} emphasize robustness, interpretability, and behavioral patterns rather than scalar returns alone. For storage control, the corresponding question is whether a controller has internalized energy-shifting behavior or instead relies on corrective safety layers to produce low evaluated cost.

\textbf{Contributions.} To address this operational validation gap, we present a morphology-aware validation framework for SAIL-SAC (Shield-Aware Intervention Learning with Soft Actor-Critic) in distribution-network microgrids. The framework integrates three components: (1) a shielded replay SAC backbone with Inventory Teacher safety filters that enforce feasibility while recording intervention magnitudes; (2) a multi-gate validation protocol (P4 BalancedGate) assessing policies across six morphological dimensions---mid-band dwelling, price-responsive directionality, infeasible-action suppression, boundary avoidance, terminal convergence, and intervention diagnostics; and (3) cross-fidelity testing from a simplified efficiency-based battery model to higher-fidelity Thevenin battery representations, assessing cross-fidelity behavior under model fidelities unavailable during training.

The intended contribution is an engineering validation protocol, not a claim of a universally superior DRL optimizer. The study therefore asks whether selected controllers produce interpretable storage trajectories under AC network settlement, and treats baseline and protocol comparisons as diagnostic evidence for the validation framework.

We evaluate the framework on an IEEE 33-bus distribution-network stress case and report full-year results for six P4-selected candidates satisfying morphological gates: mid-band dwelling $>5\%$, infeasible actions $<30\%$. For this selected set, mean objective cost is $11.788 \pm 0.120$\,M\textyen{} with $20.4 \pm 15.6$\% mid-band dwelling across full-year (365-day) cross-year evaluation (2023 training, 2024 testing) under the simple efficiency-based model. Cross-fidelity analysis shows small objective variation across battery models: 11.788\,M\textyen{} (simple EBM), 11.740\,M\textyen{} (Thevenin loss-only), 11.739\,M\textyen{} (Thevenin Rint-only), and 11.737\,M\textyen{} (Thevenin full PBM). The logged P4 evaluation records zero shield activation in the full-year six-seed evaluation but a nonzero normalized inventory-teacher action gap of $0.301 \pm 0.087$; intervention transparency is therefore a reported outcome rather than an assumed property.

\textbf{Paper organization.} Section~\ref{sec:system_model} formalizes the IEEE 33-bus microgrid testbed, battery morphology metrics, and cross-fidelity model hierarchy. Section~\ref{sec:applied_drl} details the SAIL-SAC algorithm, Inventory Teacher shield mechanism, and P4 BalancedGate validation protocol. Section~\ref{sec:case_study} presents the six-seed case study with cross-fidelity testing, intervention-diagnostic analysis, and ablation studies. Section~\ref{sec:conclusion} discusses implications, limitations, perfect-foresight oracle assumptions, and future directions for policy internalization and field-oriented evaluation.
